\reportelement{Введение}
Цель работы --- объединить FastAPI, PostgreSQL и Redis в проект Docker Compose, настроить запуск по проверкам готовности, проверить кэширование запросов к погодному API и масштабировать приложение до трёх экземпляров. Работу я выполнил в Docker Desktop на macOS в ветке \texttt{lab-4}, созданной от \texttt{lab-3}.

\begin{table}[H]
\caption{Персональные параметры лабораторной работы}
\begin{tabularx}{\textwidth}{|p{61mm}|X|}
\hline
Параметр & Значение \\\hline
ФИО & Ивановский Роман Евгеньевич \\\hline
Группа & К3320 \\\hline
Номер ИСУ & 466021 \\\hline
Порт PostgreSQL на хосте & 5430 + 21 = 5451 \\\hline
Порт ассистента на хосте & 8000 + 21 = 8021 \\\hline
Порт Redis на хосте & 6370 + 21 = 6391 \\\hline
Префикс & ivanovsky \\\hline
Учебный ключ ассистента & secret21ivanovsky \\\hline
Имя базы данных & ivanovsky21 \\\hline
\end{tabularx}
\end{table}

Порты FastAPI, PostgreSQL и Redis опубликованы на \texttt{127.0.0.1}. В сети Docker приложение обращается к PostgreSQL по адресу \texttt{db:5432}, а к Redis --- по адресу \texttt{redis:6379}. Учебный ключ и пароль PostgreSQL передаются через локальный файл \texttt{.env}, который исключён из Git.

\chapter{Запуск FastAPI и PostgreSQL в Compose}
Сначала я описал два сервиса в \texttt{docker-compose.yml}. Сервис \texttt{db} использует образ PostgreSQL 15 и именованный том \texttt{ivanovsky-pgdata} из ЛР № 3. Сервис \texttt{app} собирается из \texttt{Dockerfile.optimized}. Оба контейнера подключаются к существующей сети \texttt{ivanovsky-backend}.

Я указал параметры через \texttt{.env} и добавил в репозиторий \texttt{.env.example} с заглушками. Зависимость \texttt{depends\_on} с условием \texttt{service\_healthy} разрешает запуск FastAPI после успешной проверки PostgreSQL командой \texttt{pg\_isready}.

\begin{code}
docker compose up -d --build
docker compose ps
curl -4 -fsS http://localhost:8021/
curl -4 -fsS http://localhost:8021/health
\end{code}

Оба контейнера имели статус \texttt{healthy}. Запрос к корневому эндпоинту вернул приветствие для \texttt{ivanovsky}, а \texttt{/health} подтвердил подключение к PostgreSQL.

\reportscreenshot{images/step-1-3.png}{Базовый запуск FastAPI и PostgreSQL и проверка HTTP-запросами}{fig:stage-1}

\chapter{Redis и кэширование прогноза погоды}
Я добавил сервис \texttt{redis} на базе образа \texttt{redis:7-alpine}. Проверка готовности выполняется командой \texttt{redis-cli ping}. FastAPI ожидает готовности обоих зависимых сервисов и использует адрес Redis из переменной \texttt{REDIS\_URL}.

В обработчике \texttt{/weather/\{city\}} я привожу имя города к нижнему регистру и формирую ключ \texttt{weather:ivanovsky:\textless{}город\textgreater{}}. При наличии записи обработчик возвращает данные из Redis с полем \texttt{source: cache}. При отсутствии записи он запрашивает прогноз в Open-Meteo, сохраняет успешный ответ на 120 секунд и возвращает \texttt{source: api}.

Эндпоинт \texttt{/cache/stats} использует обход ключей через \texttt{SCAN} и показывает владельца, настроенный TTL, число записей и список городов. Redis работает как общий кэш прогнозов. Заметки по-прежнему сохраняются в PostgreSQL.

\begin{code}
docker compose exec -T redis redis-cli ping
curl -4 -fsS http://localhost:8021/cache/stats
\end{code}

\chapter{Проверка полного стека}
После перезапуска все три сервиса получили статус \texttt{healthy}. В журналах Redis и PostgreSQL сообщили о готовности в 15:44:56 UTC, а FastAPI завершил запуск в 15:45:02 UTC.

\reportscreenshot{images/step-3-2.png}{Состояние контейнеров и порядок запуска по журналам}{fig:stage-2}

Два запроса для Москвы вернули источники \texttt{api} и \texttt{cache}. В статистике появилась запись \texttt{moscow}, а Redis показал оставшееся время жизни ключа 120 секунд.

\reportscreenshot{images/step-3-3.png}{Повторный запрос погоды для Москвы и статистика кэша}{fig:stage-3}

Команда \texttt{docker network inspect} подтвердила, что все три контейнера подключены к сети \texttt{ivanovsky-backend}.

\reportscreenshot{images/step-3-4.png}{Подключение FastAPI, PostgreSQL и Redis к общей сети}{fig:stage-4}

\chapter{Масштабирование FastAPI}
Для трёх экземпляров я подключил \texttt{docker-compose.scale.yml}. Он убирает фиксированное имя контейнера FastAPI и фиксированный порт на хосте. Порт FastAPI внутри контейнера остаётся равным 8021.

\begin{code}
docker compose -f docker-compose.yml -f docker-compose.scale.yml \
  up -d --build --wait --scale app=3
\end{code}

\reportscreenshot{images/step-4-2.png}{Три экземпляра FastAPI с портами 61880, 61874 и 61877}{fig:stage-5}

Каждый экземпляр вернул успешный ответ на \texttt{/health}. Первая реплика получила прогноз для Лондона из API, а вторая и третья прочитали ту же запись из общего Redis. В статистике отобразился ключ для города \texttt{london}.

\reportscreenshot{images/step-4-3.png}{Проверка доступности трёх реплик и общего кэша для Лондона}{fig:stage-6}

\chapter{Итоговая проверка инфраструктуры}
После проверки масштабирования я вернулся к одному экземпляру FastAPI на порту 8021 и запросил прогноз для Лондона, чтобы заполнить кэш. В финальном выводе от 24.11.2026 видны три сервиса со статусом \texttt{healthy}, том \texttt{ivanovsky-pgdata} и сеть \texttt{ivanovsky-backend}.

В Redis отобразился ключ \texttt{weather:ivanovsky:london}. Эндпоинт \texttt{/cache/stats} вернул владельца \texttt{ivanovsky}, TTL 120 секунд и город \texttt{london}. На снимке также указаны время, логин и имя хоста.

\reportscreenshot{images/step-5-1.png}{Финальная проверка сервисов, сети, тома и Redis-кэша}{fig:stage-7}

Репозиторий работы, ветка lab-4.\ \url{https://github.com/st1verka/course-cont-orch-app/tree/lab-4}

\reportelement{Заключение}
Я настроил запуск FastAPI, PostgreSQL и Redis через Docker Compose. Проверки показали, что приложение запускается после готовности зависимостей, а повторный запрос погоды обслуживается из Redis. Три реплики приложения успешно подключились к PostgreSQL и прочитали общую запись кэша. После проверки я восстановил обычный запуск с одним экземпляром.

\reportelement{Приложение А. docker-compose.yml}
Первая часть конфигурации содержит настройки PostgreSQL и Redis.

\begin{code}
name: ivanovsky-lab4

services:
  db:
    image: postgres:15
    container_name: ivanovsky-postgres
    environment:
      POSTGRES_USER: ${POSTGRES_USER}
      POSTGRES_PASSWORD: ${POSTGRES_PASSWORD}
      POSTGRES_DB: ${POSTGRES_DB}
    ports:
      - "127.0.0.1:${PG_PORT}:5432"
    volumes:
      - ivanovsky-pgdata:/var/lib/postgresql/data
    networks:
      - ivanovsky-backend
    healthcheck:
      test: ["CMD-SHELL", "pg_isready -h 127.0.0.1 -U \"$${POSTGRES_USER}\" -d \"$${POSTGRES_DB}\""]
      interval: 5s
      timeout: 3s
      retries: 5

  redis:
    image: redis:7-alpine
    container_name: ivanovsky-redis
    ports:
      - "127.0.0.1:${REDIS_PORT}:6379"
    networks:
      - ivanovsky-backend
    healthcheck:
      test: ["CMD", "redis-cli", "ping"]
      interval: 5s
      timeout: 3s
      retries: 5
\end{code}

\reportelement{Приложение А. docker-compose.yml, продолжение}
Вторая часть содержит настройки FastAPI, именованного тома и сети.

\begin{code}
  app:
    build:
      context: .
      dockerfile: Dockerfile.optimized
    container_name: ivanovsky-assistant
    environment:
      PORT: ${APP_PORT}
      APP_OWNER: ${APP_OWNER}
      APP_SECRET: ${APP_SECRET}
      REDIS_URL: redis://redis:6379/0
      CACHE_TTL: ${CACHE_TTL}
      DATABASE_URL: "postgresql://${POSTGRES_USER}:${POSTGRES_PASSWORD}@db:5432/${POSTGRES_DB}"
    ports:
      - "127.0.0.1:${APP_PORT}:${APP_PORT}"
    volumes:
      - ./config:/app/config:ro
    depends_on:
      db:
        condition: service_healthy
      redis:
        condition: service_healthy
    networks:
      - ivanovsky-backend

volumes:
  ivanovsky-pgdata:
    external: true

networks:
  ivanovsky-backend:
    external: true
\end{code}

\reportelement{Приложение Б. Переменные и масштабирование}
Файл \texttt{.env.example} с параметрами запуска и заглушками для ключа и пароля.

\begin{code}
APP_OWNER=ivanovsky
APP_SECRET=CHANGE_ME
APP_PORT=8021
POSTGRES_USER=user_21
POSTGRES_PASSWORD=CHANGE_ME
POSTGRES_DB=ivanovsky21
PG_PORT=5451
REDIS_PORT=6391
CACHE_TTL=120
\end{code}

Файл \texttt{docker-compose.scale.yml} для запуска нескольких реплик.

\begin{code}
services:
  app:
    container_name: !reset null
    ports: !override
      - target: ${APP_PORT}
        host_ip: 127.0.0.1
\end{code}

\begin{thebibliography}{9}
\bibitem{source1} Курс «Контейнеризация и оркестрация приложений». Задание лабораторной работы № 4.\ \url{https://github.com/itmo-faculty-applied-informatics/course-cont-orch-app/blob/main/labs/04-compose.md}
\bibitem{source2} Репозиторий с результатами лабораторной работы, ветка lab-4.\ \url{https://github.com/st1verka/course-cont-orch-app/tree/lab-4}
\end{thebibliography}
